\documentclass[12pt]{article}
\usepackage[utf8]{inputenc}
\usepackage[T1]{fontenc}
\usepackage{hyperref}
\title{Mbariman-Gudhinma Language Resources for K-12 Curriculum Development}
\author{Compiled by Claude Sonnet 5.5 (claude-sonnet-5-5) for Andr\'as Kornai}
\date{}
\begin{document}
\maketitle

\section{Introduction}
Mbariman-Gudhinma (ISO 639-3 zmv; Glottocode mbar1253) is an Aboriginal language, or dialect cluster, of Cape York Peninsula in northern Queensland. It belongs to the Pama-Nyungan family (Paman subgroup). It is also recorded as Rimanggudinhma and Parimankutinma. Sources describe it as extinct, probably by 2003, and associate its dialects with the Lamalama people. The Endangered Languages Catalogue still lists it, noting that people continue to identify with the language. Documentation is very thin.

\section{Primary Educational Resources}
No primers, textbooks or school materials were found in the searches. A heritage curriculum would have to be built from the sketch grammar and wordlists listed below, ideally with Lamalama community involvement.

\section{Linguistic and Reference Materials}
\begin{itemize}
\item Wikipedia overview, with names, classification and the reference to Godman's thesis: \url{https://en.wikipedia.org/wiki/Mbariman-Gudhinma_language}
\item Endangered Languages Catalogue entry: \url{https://endangeredlanguages.com/lang/3132}
\item Heritage and Culture (hgcosmos) summary sheet: \url{https://www.hgcosmos.org/assets/culture/AZ/M/Mbariman-Gudhinma.pdf}
\item The search results mention Ian Godman's 1993 BA (Hons) thesis, a sketch grammar of Rimanggudinhma (University of Queensland). No direct URL for it was found.
\end{itemize}

\section{Religious and Community Resources}
None found.

\section{Digital and Audio Resources}
A Glosbe page for zmv appeared in the results, but it appears to be an automatically generated word-list page and its quality is unverified: \url{https://glosbe.com/zmv/lib}. No audio was found.

\section{Curriculum Development Context}
The language is extinct, and no teaching programme was found. Any revitalization work would be reconstruction-based and would need to be led by Lamalama community members and linked to the regional language centre. Only a small amount of material exists, so a curriculum can use the language only for greetings, names and place terms, if the community agrees.
\end{document}